\section{Discussion and limitations}
\label{sec:scope}
The study answers the reliability question through four distinctions: the joint target event, the population over which coverage averages, the score's domain and a common geometric measure with explicit numerical resolution. The controlled area reversal supports mechanism-dependent choice rather than a universal geometry. In these experiments ILR is more area-efficient than L2 for the interior mechanism; the ordering reverses near the boundary and under the specified Dirichlet generator. TV/L2 differences remain unresolved under the declared rule. These comparisons are conditional on the common fitted log-ratio center, not score-specific optimal predictors.

For a marginal target under exchangeability, a globally calibrated joint score supplies the stated guarantee. When a scientifically relevant observable group has different variability, our core results favor auditing it explicitly; aggregate coverage can be misleading. Training-only scaling substantially improves this minority but does not grant exact group validity. Fixed-group calibration is appropriate when within-group exchangeability is credible and calibration support is adequate; the scarce-group full-domain outcomes quantify its possible information cost. Its prevalence-shift protection requires invariant within-group laws, not arbitrary distribution changes.

Original-zero membership must be defined rather than silently dropping responses. Smoothing produces a valid scalar score on original candidates but changes geometry, and full-dimensional area cannot describe uncertainty along response faces. These slices depend on the response and have no stated conditional guarantee. Finite integration remains imperfect:159 small range warnings are retained, empirical allowances are nonuniform,76 area contrasts remain unresolved, and no K10 area ordering is attempted.

The target is an observed composition, not its conditional mean, a class label or latent probabilities. All data are synthetic; five mechanisms, two dimensions and one shared ridge predictor do not establish deployed-classifier, topic-model or real-data performance. The study compares specified split scores and fixed groups, not all conformal families or arbitrary subgroups. Pointwise intervals across correlated comparisons are descriptive rather than familywise tests; minor nominal shortfalls remain visible. Group reliability is a statistical diagnostic, not proof of fairness or downstream benefit. Controlled synthetic evidence makes these evaluation failures checkable, while application-specific assessment remains necessary.
